\begin{resumo}
Modelos de Linguagem de Grande Escala (LLMs) têm sido cada vez mais utilizados no desenvolvimento de software, inclusive em tarefas de geração, explicação e correção de código. No entanto, seu uso em segurança ainda exige cautela, pois correções geradas automaticamente podem ser incompletas, inseguras ou comprometer o funcionamento original da aplicação. Este trabalho investiga a integração entre LLMs, ferramentas de Análise Estática de Segurança (SAST), testes automatizados e estratégias de contextualização para apoiar a correção de \textit{issues} de segurança em código-fonte. A metodologia proposta prevê a aplicação de um \textit{pipeline} automatizado em um projeto \textit{open source}, utilizando as vulnerabilidades efetivamente detectadas pela ferramenta SAST. Como etapa inicial, foi desenvolvida uma prova de conceito para validar o fluxo de identificação, geração de correções, reanálise estática, execução de testes e consolidação de métricas. O trabalho estabelece, assim, as bases técnicas e metodológicas para avaliar a eficácia do \textit{pipeline} no TCC 2.

\textbf{Palavras-chave}: modelos de linguagem de grande escala; análise estática de segurança; correção de vulnerabilidades; engenharia de prompt; RAG; \textit{pipeline}.
\end{resumo}
